\apendice{Anexo de sostenibilización curricular}

\section{Introducción}

En este anexo se presenta una reflexión sobre los aspectos de sostenibilidad relacionados con el desarrollo del Trabajo de Fin de Grado. En particular, se analizan las implicaciones del proyecto desde el punto de vista tecnológico, educativo, social y económico, en línea con las competencias de sostenibilidad promovidas por la CRUE.

El sistema desarrollado, basado en la recopilación, análisis y visualización de datos durante el uso de Blender, permite abordar la sostenibilidad desde una perspectiva digital centrada en la eficiencia, la optimización de procesos y la mejora del aprendizaje.

Además, el proyecto integra herramientas de software libre, analíticas de aprendizaje y visualización de información, favoreciendo el desarrollo de soluciones tecnológicas accesibles, reutilizables y sostenibles.

\section{Sostenibilidad tecnológica}

Desde el punto de vista tecnológico, el proyecto promueve la eficiencia en el uso de recursos mediante distintas decisiones de diseño orientadas a minimizar el impacto computacional del sistema.

El addon de recopilación registra únicamente cambios relevantes en la interacción del usuario, evitando almacenar información redundante. Para ello, se implementaron mecanismos de logging basado en cambios y control de frecuencia de captura, reduciendo significativamente el volumen de datos generado y el coste asociado al almacenamiento y procesamiento.

Asimismo, el uso de archivos CSV como mecanismo de persistencia contribuye a mantener una arquitectura ligera y fácilmente reutilizable, evitando infraestructuras complejas o sistemas de almacenamiento de alto consumo.

El segundo addon, encargado del análisis y visualización de datos, también incorpora medidas orientadas a la eficiencia. El procesamiento de métricas se realiza de forma diferida, evitando sobrecargar el sistema durante la captura de información.

Otro aspecto relevante es el uso de software libre como Blender y Python, lo que favorece la sostenibilidad tecnológica al promover herramientas accesibles, abiertas y mantenibles a largo plazo.

\section{Sostenibilidad educativa}

El proyecto presenta una aplicación especialmente relevante en el ámbito educativo.

El sistema desarrollado permite estudiar el proceso de aprendizaje del usuario durante el modelado 3D, proporcionando información no solo sobre el resultado final, sino también sobre cómo se construye dicho resultado.

Mediante la recopilación de datos de interacción, es posible analizar aspectos como:
\begin{itemize}
    \item Estrategias de modelado.
    \item Ritmo de trabajo.
    \item Uso de herramientas.
    \item Evolución temporal del aprendizaje.
\end{itemize}

Este enfoque favorece una educación más personalizada, permitiendo comprender mejor las necesidades reales del alumnado y adaptar los procesos de enseñanza.

La actividad desarrollada en la Escuela de Arte y Superior de Diseño de Burgos permitió validar el sistema en un entorno educativo real y mostrar cómo la informática puede integrarse en disciplinas creativas como la animación y el diseño de producto.

Además, el uso de herramientas analíticas dentro de Blender contribuye a acercar conceptos relacionados con el desarrollo software, la visualización de datos y las analíticas de aprendizaje.

\section{Sostenibilidad social}

Desde el punto de vista social, el uso de software libre garantiza la accesibilidad del sistema a un amplio número de usuarios sin barreras económicas.

Blender y Python son herramientas gratuitas y ampliamente utilizadas en contextos educativos y profesionales, lo que facilita la reutilización del sistema desarrollado en distintos entornos.

La modularidad del sistema permite además que otros desarrolladores puedan ampliar o adaptar el proyecto a nuevas necesidades, favoreciendo la transferencia de conocimiento y la colaboración abierta.

Otro aspecto importante es la gestión ética de los datos recopilados. El sistema incorpora mecanismos explícitos de consentimiento del usuario antes de iniciar la captura de información, garantizando un uso responsable de los datos y respetando la privacidad del usuario.

\section{Sostenibilidad económica}

El proyecto presenta ventajas desde el punto de vista económico debido al uso de herramientas de software libre y a la ausencia de infraestructuras complejas de almacenamiento o procesamiento.

La arquitectura basada en archivos CSV y addons independientes permite desarrollar soluciones funcionales con un coste reducido, facilitando además su reutilización en futuros proyectos educativos o de investigación.

Esta reutilización incrementa el valor del sistema a largo plazo y reduce costes de desarrollo futuros.

\section{Reflexión personal}

El desarrollo de este Trabajo de Fin de Grado ha permitido aplicar competencias relacionadas con la sostenibilidad desde distintas perspectivas.

A nivel tecnológico, el proyecto ha puesto de manifiesto la importancia de diseñar sistemas eficientes y modulares, capaces de minimizar el consumo de recursos y facilitar el mantenimiento futuro.

En el ámbito educativo, el trabajo ha permitido comprender cómo las analíticas de aprendizaje pueden utilizarse para mejorar la comprensión de los procesos formativos, especialmente en entornos creativos y tridimensionales.

Asimismo, el uso de software libre ha reforzado la importancia de desarrollar herramientas accesibles y reutilizables, favoreciendo la democratización del conocimiento y el acceso a la tecnología.

\section{Conclusión}

En conjunto, el proyecto contribuye a la sostenibilidad desde múltiples perspectivas.

Desde el punto de vista tecnológico, promueve sistemas eficientes, ligeros y basados en software libre. En el ámbito educativo, facilita el análisis del proceso de aprendizaje y favorece metodologías más centradas en el usuario. A nivel social, garantiza accesibilidad y respeto por la privacidad de los usuarios.

Estas características reflejan la integración de competencias de sostenibilidad dentro del desarrollo del Trabajo de Fin de Grado, mostrando cómo la tecnología puede utilizarse de forma responsable para generar herramientas útiles, accesibles y sostenibles.

Más información en el documento de la CRUE \url{https://www.crue.org/wp-content/uploads/2020/02/Directrices_Sosteniblidad_Crue2012.pdf}.

% Pendiente revisar longitud final del anexo (objetivo: 600-800 palabras)

